% Fokus Nullstellen – Einheit 4 (quadratische-funktionen.md), Hauptnummern 6 bis 18
\setcounter{aufgabe}{5}
\einheitenkopf{Einheit 4 von 4 · Nullstellen und Schnittpunkte berechnen}
\zweigzeile{Hier lernst du, die Nullstellen einer Parabel und ihre Schnittpunkte mit einer Geraden zu berechnen · ab Kl. 10 (am Gymnasium ab Kl. 9) · P10 oft · baut auf: Scheitelpunktform (Einheit 2), Normalform (Einheit 3), Nullstelle einer Geraden, binomische Formeln, Quadratwurzeln, quadratische Gleichungen}

\begin{aufgabe}{Ich erkenne, welche Gleichung ich lösen muss. Schreibe bei a) bis c) die Gleichung auf, deren Lösungen die Nullstellen von $f$ sind. Rechne nichts aus.}
\begin{gleichungsraster}[1]
\gl{f(x) = (x - 5)^2 - 4} & \gl{f(x) = x^2 + 7x - 8} \\[7mm]
\gl{f(x) = -2x^2 + 8} & \\
\end{gleichungsraster}
\vspace{7mm}
Schreibe bei d) und e) die Gleichung auf, deren Lösungen die x-Werte der Schnittpunkte von Parabel und Gerade sind. Rechne nichts aus.
\begin{gleichungsraster}[1]
\gl{f(x) = x^2 - 1;\quad g(x) = x + 5} & \gl{f(x) = (x + 1)^2;\quad g(x) = -x + 3} \\
\end{gleichungsraster}
\end{aufgabe}

\begin{aufgabe}{Ich kann die Nullstellen einer Parabel in Scheitelpunktform berechnen. Setze $f(x) = 0$ und löse die Gleichung durch Wurzelziehen. Gib beide Nullstellen an.}
\beispiel{f(x) &= (x - 1)^2 - 9 \\ (x - 1)^2 - 9 &= 0 &&\mid +9 \\ (x - 1)^2 &= 9 &&\mid \sqrt{\phantom{x}} \\ x - 1 &= 3 \quad\text{oder}\quad x - 1 = -3 &&\mid +1 \\ x_1 &= 4 \quad\text{und}\quad x_2 = -2}
\begin{gleichungsraster}[4]
\gl{f(x) = x^2 - 25} & \gl{f(x) = (x - 4)^2 - 1} \\[7mm]
\gl{f(x) = (x - 6)^2 - 4} & \gl{f(x) = (x - 3)^2 - 25} \\
\end{gleichungsraster}
\end{aufgabe}

\begin{aufgabe}{Ich kann die Nullstellen einer Parabel in Scheitelpunktform berechnen – weiter. Berechne die Nullstellen. Gib alle Nullstellen an, die es gibt.}
\begin{gleichungsraster}[4]
\gl{f(x) = (x + 2)^2 - 9} & \gl{f(x) = (x + 5)^2 - 4} \\[7mm]
\gl{f(x) = (x - 7)^2} & \gl{f(x) = (x + 1)^2 + 4} \\
\end{gleichungsraster}
\end{aufgabe}

\begin{aufgabe}{Ich kann ohne Rechnung sagen, wie viele Nullstellen eine Parabel hat. Sieh dir bei a) bis e) den Scheitel und die Öffnung an und kreuze an, wie viele Nullstellen die Parabel hat. Rechne nicht.}
\begin{geruest}
\gz{a}{$f(x) = (x - 3)^2 - 2$}{\kreuz{keine}\kreuz{eine}\kreuz{zwei}}
\gz{b}{$f(x) = (x + 4)^2 + 1$}{\kreuz{keine}\kreuz{eine}\kreuz{zwei}}
\gz{c}{$f(x) = (x + 6)^2$}{\kreuz{keine}\kreuz{eine}\kreuz{zwei}}
\gz{d}{$f(x) = -(x - 2)^2 + 5$}{\kreuz{keine}\kreuz{eine}\kreuz{zwei}}
\gz{e}{$f(x) = -(x + 1)^2 - 3$}{\kreuz{keine}\kreuz{eine}\kreuz{zwei}}
\end{geruest}
\begin{teile}
\teil Begründe, warum du für die Nullstellen $f(x) = 0$ setzt.
\teil Jana sagt: „$f(x) = -(x - 4)^2 - 1$ hat zwei Nullstellen, weil der Scheitel unter der x-Achse liegt.“ Hat Jana recht? Begründe.
\end{teile}
\end{aufgabe}

\begin{aufgabe}{Ich kann die Nullstellen einer Parabel in Normalform berechnen. Setze $f(x) = 0$ und löse die Gleichung mit der p-q-Formel.}
\begin{teile}
\teil Ergänze die angefangene Rechnung für $f(x) = x^2 - 2x - 15$.
\end{teile}
\rechnung{x^2 - 2x - 15 &= 0 && p = -2,\ q = -15 \\ x_{1,2} &= -\frac{-2}{2} \pm \sqrt{\left(\frac{-2}{2}\right)^2 - (-15)} \\ x_{1,2} &= 1 \pm \sqrt{1 + 15} \\ x_{1,2} &= 1 \pm \text{\leerfeld} \\ x_1 &= \text{\leerfeld} \qquad x_2 = \text{\leerfeld}}
\medskip
Berechne die Nullstellen.
\begin{gleichungsraster}[4]
\gl{f(x) = x^2 - 8x + 12} & \gl{f(x) = x^2 + 4x - 21} \\[7mm]
\gl{f(x) = x^2 + 10x + 16} & \\
\end{gleichungsraster}
\end{aufgabe}

\begin{aufgabe}{Ich kann die Nullstellen einer Parabel in Normalform berechnen – weiter. Berechne die Nullstellen. Gib alle Nullstellen an, die es gibt.}
\begin{gleichungsraster}[4]
\gl{f(x) = x^2 - 5x + 6} & \gl{f(x) = x^2 + 3x - 10} \\[7mm]
\gl{f(x) = x^2 - 10x + 25} & \gl{f(x) = x^2 + 4x + 7} \\
\end{gleichungsraster}
\end{aufgabe}

\begin{aufgabe}{Ich kann Nullstellen berechnen, wenn die Wurzel nicht aufgeht. Berechne die Nullstellen. Gib sie genau mit der Wurzel an und zusätzlich auf zwei Nachkommastellen gerundet.}
\begin{gleichungsraster}[4]
\gl{f(x) = (x - 1)^2 - 3} & \gl{f(x) = (x + 3)^2 - 7} \\[7mm]
\gl[5]{f(x) = x^2 - 4x - 2} & \\
\end{gleichungsraster}
\end{aufgabe}

\begin{aufgabe}{Ich kann Nullstellen berechnen, wenn vor $x^2$ eine Zahl steht. Teile die Gleichung zuerst durch die Zahl vor $x^2$ und löse sie dann mit der p-q-Formel.}
\begin{teile}
\teil Ergänze die angefangene Rechnung für $f(x) = 2x^2 + 4x - 16$.
\end{teile}
\rechnung{2x^2 + 4x - 16 &= 0 &&\mid :2 \\ x^2 + 2x - 8 &= 0 && p = 2,\ q = -8 \\ x_{1,2} &= -1 \pm \sqrt{1 + 8} \\ x_{1,2} &= -1 \pm \text{\leerfeld} \\ x_1 &= \text{\leerfeld} \qquad x_2 = \text{\leerfeld}}
\medskip
Berechne die Nullstellen.
\begin{gleichungsraster}[5]
\gl{f(x) = 3x^2 - 21x + 30} & \gl{f(x) = -x^2 + 4x + 5} \\
\end{gleichungsraster}
\end{aufgabe}

\begin{aufgabe}{Ich kann Nullstellen berechnen, wenn vor der Klammer eine Zahl steht. Bringe zuerst die Zahl hinter der Klammer auf die andere Seite, teile dann durch die Zahl vor der Klammer und ziehe die Wurzel.}
\begin{gleichungsraster}[5]
\gl{f(x) = 2(x - 3)^2 - 18} & \gl{f(x) = -(x + 4)^2 + 1} \\
\end{gleichungsraster}
\end{aufgabe}

\begin{aufgabe}{Ich kann selbst entscheiden, wie ich die Nullstellen berechne. Berechne die Nullstellen. Bleibt eine Wurzel stehen, gib die Nullstellen genau und auf zwei Nachkommastellen gerundet an.}
\begin{gleichungsraster}[5]
\gl{f(x) = 2x^2 - 2x - 24 \quad \text{(P10 2020)}} & \gl{f(x) = x^2 + 4x - 1 \quad \text{(P10 2025)}} \\
\end{gleichungsraster}
\vspace{7mm}
Zeige bei c) durch Ausmultiplizieren, dass die Gleichung stimmt. Berechne dann die Nullstellen von $f(x) = x^2 - 8x + 9$. (P10 2017)
\begin{gleichungsraster}[6]
\gl{(x - 4)^2 - 7 = x^2 - 8x + 9} & \\
\end{gleichungsraster}
\end{aufgabe}

\begin{aufgabe}{Ich finde den Fehler beim Berechnen von Nullstellen. Lina berechnet die Nullstellen von $f(x) = (x + 1)^2 - 49$ so:}
\rechnung{(x + 1)^2 - 49 &= 0 &&\mid +49 \\ (x + 1)^2 &= 49 &&\mid \sqrt{\phantom{x}} \\ x + 1 &= 7 &&\mid -1 \\ x &= 6}
\begin{teile}
\teil Finde den Fehler, benenne ihn und korrigiere die Rechnung.
\end{teile}
Berechne die Nullstellen.
\begin{gleichungsraster}[4]
\gl{f(x) = (x - 5)^2 - 36} & \\
\end{gleichungsraster}
\end{aufgabe}

\begin{aufgabe}{Ich finde den Fehler beim Berechnen von Nullstellen – weiter. Ben berechnet die Nullstellen von $f(x) = 2x^2 + 12x + 10$ so:}
\rechnung{2x^2 + 12x + 10 &= 0 && p = 12,\ q = 10 \\ x_{1,2} &= -6 \pm \sqrt{36 - 10} \\ x_{1,2} &= -6 \pm \sqrt{26} \\ x_1 &\approx -0{,}90 \qquad x_2 \approx -11{,}10}
\begin{teile}
\teil Finde den Fehler, benenne ihn und korrigiere die Rechnung.
\end{teile}
Berechne die Nullstellen.
\begin{gleichungsraster}[5]
\gl{f(x) = 3x^2 - 12x - 36} & \\
\end{gleichungsraster}
\end{aufgabe}

\begin{aufgabe}{Ich kann mit Nullstellen berechnen, wo ein geworfener Ball landet. Tim wirft einen Ball. Die Flugbahn wird beschrieben durch $h(x) = -0{,}1\,(x - 4)^2 + 3{,}6$. Dabei ist $x$ die waagerechte Entfernung von Tim in m und $h(x)$ die Höhe des Balls über dem Boden in m.}
\begin{teile}
\teil Berechne, in welcher Höhe der Ball Tims Hand verlässt.
\teil Berechne, in welcher Entfernung von Tim der Ball auf dem Boden aufkommt.
\teil Die Gleichung in b) hat zwei Lösungen. Begründe, warum nur eine davon zur Situation passt.
\end{teile}
\end{aufgabe}
